\subsection*{4.4.3 Putting It Together}
Here we finally incorporate Lemma16 and Lemma17.
\begin{proof}[Proof of Lemma15]
By Lemma16 and Lemma17, with probability at least $1-\frac14\cdot2^{-C\log n/\log\log n}-1/n\ge1-\frac12\cdot2^{-C\log n/\log\log n}$ over $z\sim Z$, we can write $f=f_{\mathcal A}\oplus f'_{\mathcal B}\oplus f''_{\mathcal B}$, where the three functions are over disjoint set of variables, and it holds that for each $T\in\{\mathcal A,\mathcal B,\mathcal B'\}$,
\[
|(\widetilde f_T)_z(Y')-\mathbb E[\mathfrak f_T]|
\le\frac14\cdot2^{-C\log n/\log\log n}
\]
for any $\delta$-biased distribution $Y'$. Using the XOR lemma for small-biased spaces (see Lemma6), taking into account that our distribution $Y$ is in fact $\delta^3$-biased, we conclude that
\[
|\mathbb E[\widetilde f_z(Y)]-\mathbb E[\mathfrak f]|
\le16^3\cdot2\cdot\frac14\cdot2^{-C\log n/\log\log n}
\le\frac12\cdot2^{-C\log n/\log\log n},
\]
and the lemma follows.
\end{proof}
\subsection*{4.4.4 $I^g$ Almost Always Happens}
We keep using the notations of Section4.4.2. Specifically, recall that $g=f'_{\mathcal B}=\bigoplus_{i\in[m']}g_i$ for $m'=|\mathcal C\setminus X_z|$, and for $q\in[Q]$, $B_q\subseteq[m']$ is the set of terms in $g$ that belong to the $q$-th bucket. Also, for $\vec g=(\widehat g_1,\ldots,\widehat g_{m'})$,
\[
S_{\vec k}(\vec g)=\sum_{\substack{I\subseteq[m']\\K(I)=\vec k}}\prod_{i\in I}\widehat g_i.
\]
Recall that $I^g(x)\in\{0,1\}$ is1 if and only if
\[
\sum_{\substack{\vec k\in\mathbb N^Q\\\|\vec k\|_{(c)}>A}}|S_{\vec k}(\vec g(x))|\le2^{-A/1024},
\]
where $A=\bar C\log n$ and $\|\vec k\|_{(c)}=\sum_{q\in[Q]}c^q\vec k[q]$.
\begin{proof}[Proof of Lemma22]
As in Section2, we define $R_{\vec k}(\vec g)=S_{\vec k}^2(\vec g)\cdot\prod_{q\in[Q]}\vec k[q]!$.
By Lemma3, to prove the bound on $\Pr[I^g(Y)=0]$ it is sufficient to prove that for every $\vec k\in\mathbb N^Q$ with $\|\vec k\|_{(c)}\le A$ we have that
\[
\mathbb E[R_{\vec k}(\vec g(Y))]\le2^{-\|\vec k\|_{(c)}/8}.
\]
By now a standard course of action, we aim at bounding the spectral norm of the function $R_{\vec k}(\vec g)$, together with its expectation under the uniform distribution. To this end, let us define, for $q\in[Q]$ and an integer $\ell$,
\[
\check S_{\ell,q}=\sum_{\substack{I\subseteq B_q\\|I|=\ell}}\prod_{i\in I}\widehat g_i,
\]
so $R_{\vec k}(\vec g)=\prod_{q\in[Q]}\check S_{\vec k[q],q}^2\vec k[q]!$. First, we record that:
\begin{authorclaim}{Claim25}
For every $i\in[m']$, $\mathbb E[\widehat g_i^2]\le2^{-(2-2\beta)w_i}$.
\end{authorclaim}
\begin{proof}
Let $V_i\subseteq[n]$ be the set of variables read by $g_i$, of cardinality $w_i$, and let $\ell_i=|V_i\cap\{j\in[n]:z_j=1\}|$ be the number of live variables read by $g_i$. Note that
\[
\widetilde g_i(x)=\mathbb E[g_i(x+z\wedge U)]
=\begin{cases}0&\text{if there exists }j\in V_i\text{ such that }x_j=z_j=0,\\2^{-\ell_i}&\text{otherwise}.\end{cases}
\]
Then,
\begin{align*}
\mathbb E[\widetilde g_i^2]
&=2^{-2\ell_i}\Pr_{x\sim U}[\text{for every }j\in V_i\text{ s.t. }z_j=0\text{ it holds that }x_j=1]\\
&=2^{-2\ell_i}2^{-(w_i-\ell_i)}=2^{-w_i-\ell_i}.
\end{align*}
Recalling that $\ell_i\ge(1-\beta)w_i$ ($g_i$ is good under $z$), we have $\mathbb E[\widetilde g_i^2]\le2^{-(2-\beta)w_i}$. Let $h_i=1-g_i$, and note that
\[
\mathbb E[\widehat h_i^2]=\operatorname{Var}[\widehat h_i]
=\frac{\operatorname{Var}[\widetilde g_i]}{\mathbb E[h_i]^2}
\le4\cdot\mathbb E[\widetilde g_i^2]\le4\cdot2^{-(2-\beta)w_i}
\le2^{-(2-2\beta)w_i}.
\]
The fact that $\operatorname{Var}[\widehat h_i]=\operatorname{Var}[\widehat g_i]=\mathbb E[\widehat g_i^2]$ finishes the proof.
\end{proof}
Now,
\begin{align*}
\mathbb E[\check S_{\ell,q}^2]
&=\sum_{\substack{I\subseteq B_q\\|I|=\ell}}\prod_{i\in I}\mathbb E[\widehat g_i^2]
\le\sum_{\substack{I\subseteq B_q\\|I|=\ell}}\prod_{i\in I}2^{-(2-2\beta)w_i}\\
&\le\sum_{\substack{I\subseteq B_q\\|I|=\ell}}2^{-(2-2\beta)c^q\ell}
\le\binom{m_q}{\ell}2^{-(2-2\beta)c^q\ell}\\
&\le\frac{3^{c^q\ell}e^\ell2^{-(2-2\beta)c^q\ell}}{\ell!}
\le\frac{2^{-c^q\ell/4}}{\ell!}.
\end{align*}
Plugging it in our expression for $R_{\vec k}$, we get
\begin{equation}
\mathbb E[R_{\vec k}(\vec g)]
=\prod_{q\in[Q]}\mathbb E[\check S_{\vec k[q],q}^2\vec k[q]!]
\le\prod_{q\in[Q]}2^{-c^q\vec k[q]/4}=2^{-\|\vec k\|_{(c)}/4}.\tag{13}
\end{equation}
Finally, let us bound $L_1(R_{\vec k}(\vec g))$. In Equation(12) we established the fact that $L_1(S_{\vec k}(\vec g))\le12^{\|\vec k\|_{(c)}}\le12^A$. Thus,
\begin{align*}
L_1(R_{\vec k}(\vec g))
&\le12^{2A}\prod_{q\in[Q]}\vec k[q]!
\le12^{2A}e^{\sum_{q\in[Q]}\vec k[q]\ln\vec k[q]}\\
&\le12^{2A}e^{(\ln A)\sum_{q\in[Q]}\vec k[q]}.
\end{align*}
As $\|\vec k\|_{(c)}=\sum_{q\in[Q]}c^q\vec k[q]\le A$ and $c^q\ge C\log\log n$ (see Claim18), $\sum_{q\in[Q]}\vec k[q]\le A/(C\log\log n)$ and we get
\[
L_1(R_{\vec k}(\vec g))
\le12^{2A}e^{\ln A\cdot A/(C\log\log n)}
\le12^{2A}2^{(\bar C/C)\log n}\le n^{10\bar C}.
\]
Note that $\delta\le\frac1{32}n^{-10\bar C}2^{-A/4}$. Thus, by Claim5,
\[
\mathbb E[R_{\vec k}(\vec g(Y))]
\le2^{-\|\vec k\|_{(c)}/4}+\delta\cdot n^{10\bar C}
\le2^{-\|\vec k\|_{(c)}/8},
\]
and we are done with bounding $\Pr[I^g(Y)=0]$. For the bound on $\mathbb E[S_{\vec k}^2(\vec g(Y))]$, simply observe that $\mathbb E[S_{\vec k}^2(\vec g(Y))]<\mathbb E[R_{\vec k}(\vec g(Y))]$.
\end{proof}
